\documentclass[11pt]{article}
% ============================================================================
% preamble.tex — Estilo común de los manuales de práctica SISELA (FIME / UANL)
%
% Reproduce (de forma aproximada) el estilo de los PDF originales
% (~/Descargas/manuales/Manuales/pN.pdf), de los que no existe fuente .tex.
% Compilar con:  pdflatex  (o  latexmk -pdf pN.tex)
% ============================================================================

\usepackage[utf8]{inputenc}
\usepackage[T1]{fontenc}
\usepackage[spanish,es-tabla]{babel}
\usepackage[a4paper,margin=2.5cm,headheight=15pt]{geometry}

% --- Tipografía: serif para el cuerpo, sans para los títulos ---
\usepackage{libertinus}
\usepackage[scaled=0.92]{helvet}   % sans para encabezados
\usepackage{microtype}

% --- Matemáticas ---
\usepackage{amsmath,amssymb}
\usepackage{siunitx}
\sisetup{detect-all, per-mode=symbol}

% --- Tablas y figuras ---
\usepackage{booktabs}
\usepackage{array}
\usepackage{multirow}
\usepackage{caption}
\captionsetup{font=small,labelfont=bf,labelsep=period}
\usepackage{graphicx}
\usepackage{float}

% --- Diagramas ---
\usepackage{tikz}
\usetikzlibrary{shapes.geometric,arrows.meta,positioning,calc}
\tikzset{
  block/.style   = {draw, rectangle, minimum width=2.6cm, minimum height=0.9cm,
                    align=center, font=\footnotesize},
  term/.style    = {draw, rounded corners, minimum width=2.2cm, minimum height=0.8cm,
                    align=center, font=\footnotesize, fill=black!5},
  decision/.style= {draw, diamond, aspect=2, align=center, font=\scriptsize,
                    inner sep=1pt},
  flow/.style    = {-{Latex[length=2mm]}, thick},
}

% --- Código ---
\usepackage{xcolor}
\definecolor{codebg}{gray}{0.96}
\definecolor{codekw}{rgb}{0.13,0.29,0.53}
\definecolor{codecm}{rgb}{0.40,0.45,0.40}
\usepackage{listings}
\lstdefinestyle{sisela}{
  backgroundcolor=\color{codebg}, basicstyle=\ttfamily\footnotesize,
  keywordstyle=\color{codekw}\bfseries, commentstyle=\color{codecm}\itshape,
  stringstyle=\color{codekw}, numbers=left, numberstyle=\tiny\color{gray},
  numbersep=6pt, frame=none, breaklines=true, showstringspaces=false,
  columns=fullflexible, xleftmargin=1.8em, aboveskip=1em, belowskip=1em,
  literate={á}{{\'a}}1 {é}{{\'e}}1 {í}{{\'i}}1 {ó}{{\'o}}1 {ú}{{\'u}}1
           {ñ}{{\~n}}1 {Á}{{\'A}}1 {°}{{\textdegree}}1 {µ}{{$\mu$}}1
           {✓}{{\checkmark}}1,
}
\lstset{style=sisela}
\newcommand{\code}[1]{\texttt{\small #1}}

% --- Listas ---
\usepackage{enumitem}
\setlist{nosep, leftmargin=1.6em}

% --- Encabezado / pie de página ---
\usepackage{fancyhdr}
\usepackage{lastpage}
\pagestyle{fancy}
\fancyhf{}
\fancyhead[L]{\footnotesize\sffamily Sistemas Eléctricos de Aeronaves — Práctica \practnum}
\fancyhead[R]{\footnotesize\sffamily FIME / UANL}
\fancyfoot[C]{\footnotesize P\'agina \thepage\ de \pageref{LastPage}}
\renewcommand{\headrulewidth}{0.4pt}
\renewcommand{\footrulewidth}{0.4pt}

% --- Títulos de sección en sans-serif ---
\usepackage{titlesec}
\titleformat{\section}{\normalfont\Large\sffamily\bfseries}{\thesection}{0.7em}{}
\titleformat{\subsection}{\normalfont\large\sffamily\bfseries}{\thesubsection}{0.6em}{}
\titleformat{\subsubsection}{\normalfont\normalsize\sffamily\bfseries}{\thesubsubsection}{0.5em}{}

% --- hyperref al final ---
\usepackage[hidelinks,pdfusetitle]{hyperref}

% ============================================================================
% Comandos de portada y cajas
% ============================================================================
\newcommand{\practnum}{N}          % lo redefine cada pN.tex
\newcommand{\practtitle}{}         % idem

\newcommand{\portada}{%
  \thispagestyle{empty}
  \begin{center}
    {\large Universidad Autónoma de Nuevo León}\\[2pt]
    {\normalsize Facultad de Ingeniería Mecánica y Eléctrica}\\[2.2cm]
    \rule{\linewidth}{0.4pt}\\[0.4cm]
    {\Huge\sffamily\bfseries Práctica \practnum}\\[0.5cm]
    {\large \practtitle}\\[0.3cm]
    \rule{\linewidth}{0.4pt}\\[2.5cm]
    {\normalsize Laboratorio de Sistemas Electrónicos de Aeronaves\\(SISELA)}\\[1.5cm]
    {\normalsize Ciclo Escolar 2026}\\[3cm]
  \end{center}
}

\newenvironment{resumen}{\par\noindent\rule{\linewidth}{0.4pt}\par\smallskip
  \small\itshape\noindent\textbf{Resumen Ejecutivo} ---\ }{\par\smallskip
  \noindent\rule{\linewidth}{0.4pt}\par}

\newenvironment{nota}[1][Nota]{\par\smallskip\noindent\rule{\linewidth}{0.4pt}\par\smallskip
  \small\itshape\noindent\textbf{#1} ---\ }{\par\smallskip
  \noindent\rule{\linewidth}{0.4pt}\par\smallskip}

\renewcommand{\practnum}{4}
\renewcommand{\practtitle}{Altímetro Barométrico con Sensor BMP180\\
Comunicación I2C, Compensación Digital y Análisis de Señales — Énfasis Aeronáutico}

\begin{document}
\portada

\begin{resumen}
Esta sesión implementa un altímetro barométrico funcional con el sensor digital BMP180
(Bosch). Se aborda la comunicación I2C maestro--esclavo, la lectura de 11 coeficientes de
calibración de fábrica, el algoritmo completo de compensación de temperatura y presión
según el \emph{datasheet}, y el cálculo de altitud barométrica mediante la fórmula ISA.
Se introducen conceptos de altimetría aeronáutica (QNH, QFE, QNE, altitud de transición)
y se realizan mediciones comparativas a distintas alturas. La práctica incorpora además
un módulo de \textbf{análisis de señales} sobre la serie temporal del sensor: verificación
de la frecuencia de muestreo, caracterización del ruido (desviación de Allan, resolución
efectiva), análisis espectral del piso de ruido y diseño de filtros digitales.
\end{resumen}

\tableofcontents
\newpage

% ============================================================================
\section{Introducción}

La medición precisa de la presión barométrica es uno de los pilares de la navegación
aérea. El altímetro barométrico, presente en todo panel de vuelo, determina la altitud
de la aeronave midiendo la presión atmosférica y comparándola con un modelo de atmósfera
estándar. Aun con sistemas GNSS, el altímetro barométrico sigue siendo el medio primario
de separación vertical en espacio aéreo controlado, porque proporciona una referencia de
altitud común basada en la misma presión de referencia (QNH).

El \textbf{BMP180} de Bosch Sensortec es un sensor MEMS completamente digital que se
comunica por el bus I2C. Internamente combina un puente de Wheatstone piezoresistivo, un
ADC sigma--delta de alta resolución y una memoria EEPROM con 11 coeficientes de
calibración grabados de fábrica para compensar las variaciones del proceso de fabricación.

Esta práctica profundiza en: (i)~el protocolo I2C (maestro--esclavo, direccionamiento,
registros); (ii)~la lectura e interpretación de los coeficientes de calibración;
(iii)~el algoritmo de compensación entera de 32 bits del \emph{datasheet};
(iv)~el cálculo de altitud con la fórmula ISA; (v)~la altimetría aeronáutica
(QNH/QFE/QNE); y (vi)~el \textbf{análisis de la señal de salida del sensor} como una
serie temporal muestreada (ruido, muestreo, espectro y filtrado).

% ============================================================================
\section{Objetivos de Aprendizaje}

Al finalizar esta práctica, el estudiante será capaz de:
\begin{enumerate}
  \item Explicar el principio MEMS piezoresistivo del BMP180 y su estructura interna
        (puente de Wheatstone, ADC sigma--delta, EEPROM).
  \item Configurar el bus I2C como maestro y realizar transacciones de lectura/escritura
        por dirección y registro; verificar el dispositivo mediante escaneo del bus.
  \item Leer los 11 coeficientes de calibración (AC1--AC6, B1, B2, MB, MC, MD) desde la
        EEPROM y comprender su rol en la compensación (con y sin signo).
  \item Implementar el algoritmo de compensación completo
        ($\text{UT}\to X_1\to X_2\to B_5\to T$;
         $\text{UP}\to B_6\to B_3\to B_4\to B_7\to P$) en aritmética entera de 32 bits.
  \item Calcular la altitud barométrica
        $h = \num{44330}\,\bigl(1-(P/P_0)^{1/5.255}\bigr)$ y discutir sus limitaciones.
  \item Operar con referencias altimétricas: configurar el QNH para obtener altitud real
        respecto al nivel medio del mar, diferenciando QNH, QFE y QNE.
  \item Verificar la \textbf{frecuencia de muestreo real} de la adquisición y su
        \emph{jitter}, y relacionarla con el teorema de Nyquist.
  \item Caracterizar el \textbf{ruido} de la medición de altitud (desviación estándar,
        RMS, resolución efectiva en bits, densidad espectral de potencia, desviación de
        Allan).
  \item Diseñar y evaluar \textbf{filtros digitales} (media móvil, mediana, exponencial)
        y medir su respuesta en frecuencia y su respuesta al escalón.
  \item Documentar resultados siguiendo la estructura de reporte técnico AIAA
        simplificado.
\end{enumerate}

% ============================================================================
\section{Fundamentos Teóricos}

\subsection{Presión barométrica y atmósfera estándar ISA}
La presión atmosférica al nivel del mar en condiciones ISA es
\begin{equation}
  P_0 = \SI{101325}{\pascal} = \SI{1013.25}{\hecto\pascal} = \SI{29.92}{inHg}.
\end{equation}
El modelo ISA en la troposfera ($h<\SI{11000}{\meter}$):
\begin{equation}
  P(h) = P_0\left(1 - \frac{L\,h}{T_0}\right)^{\frac{g_0}{R\,L}},
\end{equation}
con $L=\SI{0.0065}{\kelvin\per\meter}$, $T_0=\SI{288.15}{\kelvin}$,
$g_0=\SI{9.80665}{\meter\per\second\squared}$,
$R=\SI{287.05}{\joule\per\kilogram\per\kelvin}$. La ecuación inversa (fórmula
barométrica) es
\begin{equation}
  \boxed{\; h = \frac{T_0}{L}\left[1-\left(\frac{P}{P_0}\right)^{\frac{R\,L}{g_0}}\right]
        = \num{44330}\left[1-\left(\frac{P}{P_0}\right)^{1/5.255}\right]\si{\meter}. \;}
  \label{eq:isa}
\end{equation}

\begin{table}[H]\centering
\caption{Atmósfera estándar ISA: valores seleccionados.}
\begin{tabular}{rrrr}
\toprule
Altitud (m) & Temp.\ (\si{\celsius}) & Presión (hPa) & Densidad (\si{\kilogram\per\meter\cubed}) \\
\midrule
0     & 15.0  & 1013.25 & 1.2250 \\
500   & 11.8  & 954.61  & 1.1673 \\
1000  & 8.5   & 898.76  & 1.1117 \\
2000  & 2.0   & 794.95  & 1.0066 \\
3000  & $-4.5$  & 701.09  & 0.9091 \\
5000  & $-17.5$ & 540.20  & 0.7361 \\
\bottomrule
\end{tabular}
\end{table}

\begin{nota}[Gradiente de presión]
Cerca del nivel del mar la presión disminuye \SI{1}{\hecto\pascal} por cada
\SI{8.5}{\meter} de ascenso ($\approx\SI{12}{\pascal\per\meter}$). Un altímetro con
resolución de \SI{0.01}{\hecto\pascal} puede resolver diferencias de $\sim\SI{8.5}{\centi\meter}$.
\end{nota}

\subsection{Sensor BMP180: principio MEMS piezoresistivo}
\begin{table}[H]\centering
\caption{Especificaciones principales del BMP180.}
\begin{tabular}{ll}
\toprule
Parámetro & Valor \\
\midrule
Rango de presión & 300--1100 hPa ($\approx +9000$ a $-500$ m) \\
Resolución (presión) & 0.01 hPa (modo ultra-alta resolución, \code{oss}=3) \\
Precisión absoluta & $\pm 1.0$ hPa típica \\
Rango de temperatura & $-40$ a $+85$ \si{\celsius} \\
Interfaz & I2C, dirección fija 0x77 \\
Chip ID & 0x55 (registro 0xD0) \\
Alimentación & 1.8--3.6 V \\
Consumo & \SI{5}{\micro\ampere} (estándar), \SI{0.1}{\micro\ampere} (standby) \\
\bottomrule
\end{tabular}
\end{table}

\begin{table}[H]\centering
\caption{Modos de sobremuestreo (OSS) del BMP180.}
\begin{tabular}{lcccc}
\toprule
Modo & OSS & Muestras internas & Tiempo (ms) & Ruido RMS (Pa) \\
\midrule
Ultra Low Power & 0 & 1 & 4.5  & 6.0 \\
Estándar        & 1 & 2 & 7.5  & 5.0 \\
Alta resolución & 2 & 4 & 13.5 & 4.0 \\
Ultra High Res. & 3 & 8 & 25.5 & 3.0 \\
\bottomrule
\end{tabular}
\end{table}

\subsection{Protocolo I2C}
I2C usa dos líneas \emph{open-drain} con pull-up ($\approx\SI{4.7}{\kilo\ohm}$):
\textbf{SDA} (datos) y \textbf{SCL} (reloj, generado por el maestro). Cada dispositivo
tiene una dirección de 7 bits (BMP180 = $0\mathrm{x}77 = 0111\,0111_2$). El patrón de
transacción es START $\to$ dirección + R/W $\to$ ACK $\to$ datos $\to$ ACK $\to$ STOP.
La comunicación con el BMP180 sigue el patrón registro--dato:
\begin{itemize}
  \item \textbf{Escritura}: START $\to$ (\code{0x77}, W) $\to$ ACK $\to$ registro $\to$
        ACK $\to$ dato $\to$ ACK $\to$ STOP.
  \item \textbf{Lectura}: START $\to$ (\code{0x77}, W) $\to$ ACK $\to$ registro $\to$
        ACK $\to$ RESTART $\to$ (\code{0x77}, R) $\to$ ACK $\to$ datos $\to$ NACK $\to$ STOP.
\end{itemize}

\subsection{Mapa de registros y coeficientes de calibración}
\begin{table}[H]\centering
\caption{Registros principales del BMP180.}
\begin{tabular}{llll}
\toprule
Dirección & Nombre & Acceso & Descripción \\
\midrule
0xAA--0xBF & Calibración & R & 11 coeficientes (22 bytes, big-endian) \\
0xD0       & Chip ID     & R & Debe ser 0x55 \\
0xE0       & Soft Reset  & W & Escribir 0xB6 \\
0xF4       & Control     & R/W & Inicia medición \\
0xF6--0xF8 & Data        & R & Resultado (16 o 19 bits) \\
\bottomrule
\end{tabular}
\end{table}
Para iniciar una medición se escribe en 0xF4: temperatura $\to$ \code{0x2E}, esperar
\SI{4.5}{\milli\second}, leer 16 bits $\to$ UT. Presión $\to$
\code{0x34 + (oss} $\ll$ \code{6)}, esperar según el modo, leer 24 bits, desplazar
$\gg(8-\text{oss})$ $\to$ UP.

\begin{table}[H]\centering
\caption{Coeficientes de calibración del BMP180.}
\begin{tabular}{llll}
\toprule
Coef.\ & Dirección & Tipo & Rol \\
\midrule
AC1 & 0xAA & int16  & Compensación principal \\
AC2 & 0xAC & int16  & Compensación cruzada \\
AC3 & 0xAE & int16  & Compensación cruzada \\
AC4 & 0xB0 & uint16 & Factor de ganancia \\
AC5 & 0xB2 & uint16 & Factor de conversión de temp. \\
AC6 & 0xB4 & uint16 & Offset de temperatura \\
B1  & 0xB6 & int16  & Compensación de presión \\
B2  & 0xB8 & int16  & Compensación cruzada T--P \\
MB  & 0xBA & int16  & Reservado \\
MC  & 0xBC & int16  & Offset de temperatura \\
MD  & 0xBE & int16  & Factor de conversión de temp. \\
\bottomrule
\end{tabular}
\end{table}

\begin{nota}[Tipos con/sin signo]
AC4, AC5 y AC6 son \textbf{unsigned} de 16 bits (0--65535). El resto son \textbf{signed}
de 16 bits ($-32768$ a $+32767$). En MicroPython, los valores leídos del bus I2C son
siempre unsigned; se requiere conversión explícita a signed para los coeficientes que lo
necesitan: \code{val - 65536 if val >= 32768 else val}.
\end{nota}

\subsection{Algoritmo de compensación de temperatura}
\begin{align}
  X_1 &= \frac{(\text{UT} - \text{AC6})\times\text{AC5}}{2^{15}} \\
  X_2 &= \frac{\text{MC}\times 2^{11}}{X_1 + \text{MD}} \\
  B_5 &= X_1 + X_2 \\
  T   &= \frac{B_5 + 8}{2^4} \quad (\text{en \SI{0.1}{\celsius}}).
\end{align}

\subsection{Algoritmo de compensación de presión}
\begin{align}
  B_6 &= B_5 - 4000, &
  X_1 &= \frac{B_2\lfloor B_6^2/2^{12}\rfloor}{2^{11}}, &
  X_2 &= \frac{\text{AC2}\,B_6}{2^{11}}, \\
  X_3 &= X_1 + X_2, &
  B_3 &= \frac{(\text{AC1}\cdot 4 + X_3)\ll\text{oss} + 2}{4}, & & \\
  X_1 &= \frac{\text{AC3}\,B_6}{2^{13}}, &
  X_2 &= \frac{B_1\lfloor B_6^2/2^{12}\rfloor}{2^{16}}, &
  X_3 &= \frac{(X_1+X_2)+2}{4}, \\
  B_4 &= \frac{\text{AC4}\,(X_3+32768)}{2^{15}}, &
  B_7 &= (\text{UP}-B_3)(50000\gg\text{oss}). & &
\end{align}
Selección de escala: $P = B_7\cdot 2 / B_4$ si $B_7 < \mathrm{0x80000000}$, en caso
contrario $P = (B_7/B_4)\cdot 2$. Corrección final:
\begin{equation}
  P \mathrel{+}= \frac{\lfloor P/2^8\rfloor^2\cdot 3038/2^{16} \;-\; 7357\,P/2^{16} \;+\; 3791}{2^4}
  \quad(\si{\pascal}).
\end{equation}

\begin{nota}[Aritmética entera]
El \emph{datasheet} de Bosch especifica el algoritmo en aritmética entera de 32 bits. Las
divisiones por potencias de 2 se implementan como desplazamientos ($\gg n$). Esto
garantiza un resultado idéntico en cualquier plataforma (MicroPython, C++, FPGA) sin
errores de punto flotante.
\end{nota}

\subsection{Altimetría aeronáutica: QNH, QFE, QNE}
\begin{table}[H]\centering
\caption{Referencias altimétricas en aviación.}
\begin{tabular}{lll}
\toprule
Referencia & $P_0$ ajustada a & Altitud indicada \\
\midrule
QNH & Nivel medio del mar (local) & Altitud MSL \\
QFE & Nivel del aeródromo & Altura AGL \\
QNE & 1013.25 hPa (estándar) & Altitud de presión / Nivel de vuelo \\
\bottomrule
\end{tabular}
\end{table}
La ecuación~\eqref{eq:isa} se emplea con $P_0 = \text{QNH}$ para obtener la altitud real
respecto al nivel del mar. Si se conoce la elevación $h_\text{conocida}$ del punto de
medición, el QNH equivalente es
$P_0 = P / \bigl(1 - h_\text{conocida}/\num{44330}\bigr)^{5.255}$.

\subsection{ARINC 429: Label 203 (Baro-Corrected Altitude)}
En aviónica avanzada la altitud barométrica corregida se transmite en formato ARINC 429
bajo la Label 203 (octal). La palabra de 32 bits contiene: bits 1--8 label (203 octal,
orden inverso), 9--10 SDI, 11--29 dato BNR (resolución $\sim\SI{1}{ft}$), 30--31 SSM,
32 paridad impar.

% ============================================================================
\section{Consideraciones de Seguridad y Buenas Prácticas}
\begin{enumerate}
  \item \textbf{Alimentación del sensor}: el BMP180 opera a 1.8--3.6 V. El módulo GY-68
        incluye regulador LDO y acepta 3.3 V o 5 V en VCC.
  \item \textbf{Protección de pines I2C}: SDA/SCL son \emph{open-drain} a 3.3 V. No
        conectar señales de 5 V sin conversión de nivel.
  \item \textbf{Pull-ups}: el módulo GY-68 ya incluye pull-ups de \SI{4.7}{\kilo\ohm}. No
        añadir pull-ups adicionales en la protoboard.
  \item \textbf{ESD}: el BMP180 es sensible a descargas electrostáticas.
  \item \textbf{Montaje}: un mal contacto en SDA/SCL causa errores intermitentes
        difíciles de diagnosticar.
  \item \textbf{Interpretación de altitud}: las lecturas barométricas \emph{no} son
        precisión GPS. No usar para navegación real.
\end{enumerate}

% ============================================================================
\section{Equipo y Materiales}
\begin{itemize}
  \item Módulo sensor BMP180 (GY-68) con regulador y pull-ups integrados.
  \item Tarjeta de desarrollo: ESP32 DevKit o Raspberry Pi Pico (RP2040).
  \item Protoboard y 4 cables Dupont (VCC, GND, SDA, SCL).
  \item Cable USB para programación y comunicación serie (115200 baud).
  \item Osciloscopio de 2 canales (para observar el bus I2C).
  \item Analizador lógico (opcional, decodificación de protocolo I2C).
  \item PC con Python 3.10+ y la toolkit \code{tools/sisela\_signal/}
        (\code{numpy}, \code{scipy}, \code{matplotlib}, \code{pyserial}).
\end{itemize}

% ============================================================================
\section{Etapa de Diseño}

\subsection{Asignación de pines por plataforma}
\begin{table}[H]\centering
\caption{Pines I2C para cada plataforma.}
\begin{tabular}{llll}
\toprule
Señal & ESP32 (MicroPython) & RP2040 (MicroPython) & C++ (Arduino) \\
\midrule
SDA & GPIO21 & GP0 & 21 (ESP32) / 4 (Pico) \\
SCL & GPIO22 & GP1 & 22 (ESP32) / 5 (Pico) \\
VCC & 3V3 & 3V3(OUT) & 3V3 \\
GND & GND & GND & GND \\
\bottomrule
\end{tabular}
\end{table}

\subsection{Cálculos de verificación (datos de ejemplo Bosch)}
Con AC1 = 408, AC2 = $-72$, AC3 = $-14383$, AC4 = 32741, AC5 = 32757, AC6 = 23153,
B1 = 6190, B2 = 4, MB = $-32768$, MC = $-8711$, MD = 2868, UT = 27898, UP = 23843:
\begin{align*}
  X_1 &= (27898 - 23153)\times 32757/2^{15} = 4743, \\
  X_2 &= (-8711)\times 2^{11}/(4743 + 2868) = -2344, \\
  B_5 &= 4743 + (-2344) = 2399, \\
  T   &= (2399 + 8)/16 = 150 \;\Rightarrow\; \SI{15.0}{\celsius}.
\end{align*}
Si $P = \SI{69964}{\pascal}$ y $P_0 = \SI{101325}{\pascal}$:
$h = \num{44330}\,(1-(69964/101325)^{1/5.255}) \approx \SI{3139}{\meter}$.

\subsection{Diagrama general de fases}
\begin{figure}[H]\centering
\begin{tikzpicture}[node distance=0.9cm]
  \node[term] (a) {Iniciar P4};
  \node[block, below=of a] (b) {Escaneo I2C\\Verificar 0x77};
  \node[block, below=of b] (c) {Leer 11 coef.\\de EEPROM};
  \node[decision, below=of c] (d) {Menú de\\modos (1--7)};
  \node[block, below=of d] (e) {Casos 1--7\\(ver \S7)};
  \draw[flow] (a)--(b); \draw[flow] (b)--(c); \draw[flow] (c)--(d); \draw[flow] (d)--(e);
\end{tikzpicture}
\caption{Diagrama de flujo general de la Práctica 4.}
\end{figure}

% ============================================================================
\section{Procedimiento Experimental}

\subsection{Paso 1: preparación del entorno}
\begin{enumerate}
  \item Conectar el módulo BMP180 según la tabla de pines de la plataforma.
  \item Verificar con multímetro 3.3 V en VCC del módulo.
  \item Cargar el programa (incluir \code{lib/bmp180.py} y \code{lib/siglab.py} para
        MicroPython).
  \item Abrir terminal serie a 115200 baud.
  \item Verificar que el escaneo I2C detecte el dispositivo en 0x77.
\end{enumerate}

\subsection{Caso 1: datos crudos y coeficientes de calibración}
El programa lee y muestra los 11 coeficientes almacenados en la EEPROM, seguido de
lecturas continuas de UT y UP. Anotar los valores de los 11 coeficientes en la
Tabla~\ref{tab:coef}.

\subsection{Caso 2: temperatura y presión compensadas}
El modo muestra el resultado del algoritmo paso a paso ($X_1$, $X_2$, $B_5$, $B_6$,
$B_3$, $B_4$, $B_7$). Comparar los valores intermedios con los cálculos manuales de la
\S6.2.

\subsection{Caso 3: altímetro barométrico con ajuste QNH}
El modo implementa un altímetro que muestra altitud en metros y pies con ajuste de QNH
(comando \code{q}) o calibración por altitud conocida (\code{aXXX}). Si la elevación del
edificio de FIME es $\approx\SI{540}{\meter}$, calcular el QNH equivalente.

\subsection{Caso 4: monitor CSV continuo}
Emite \code{timestamp\_ms,temp\_C,pressure\_hPa,altitude\_m} a 5 Hz. Alimenta directamente
la herramienta \code{altimeter\_gui.py}.

\subsection{Caso 5: comparativa de alturas}
Guía al estudiante para tomar mediciones en $\ge 3$ niveles distintos, calculando el
promedio de 10 muestras/punto, la desviación estándar $\sigma_h$, y las diferencias de
altura $\Delta h$ entre puntos consecutivos y respecto al primero.

\subsection{Caso 6: análisis de ruido y muestreo}
\begin{enumerate}
  \item Con el sensor quieto, capturar un bloque de $N$ muestras de altitud a una
        frecuencia $F_s$ seleccionable (el firmware usa \code{siglab.BlockSampler}).
  \item El reporte indica: $F_s$ real vs solicitada, \emph{jitter} de muestreo $\sigma(\Delta t)$,
        desviación estándar de la altitud $\sigma_h$, RMS del ruido, y resolución
        efectiva en bits.
  \item Repetir con \code{oss} = 0, 1, 2, 3.
  \item Guardar el CSV y en la PC ejecutar:
\begin{lstlisting}[language=bash]
python -m sisela_signal spectrum     --file cap.csv --col alt_m --psd
python -m sisela_signal characterize --file cap.csv --col alt_m --mode allan
\end{lstlisting}
\end{enumerate}

\subsection{Caso 7: filtro digital en vivo}
El modo aplica en la placa un filtro (media móvil, mediana o EMA) y transmite
\code{t\_us,alt\_raw,alt\_filt}. Dejar el sensor quieto 20 s y luego subir un piso de
escaleras (escalón). En la PC:
\begin{lstlisting}[language=bash]
python -m sisela_signal filter       --file cap.csv --col alt_raw --kind movavg --n 8 --bode
python -m sisela_signal characterize --file cap.csv --col alt_filt --mode step
\end{lstlisting}

% ============================================================================
\section{Análisis de Señales}
\label{sec:senales}

El BMP180 es un sensor \emph{digital}: no hay una señal analógica que inyectar con el
generador de funciones. El análisis de señales de esta práctica se realiza sobre la
\textbf{serie temporal digital} $y[n]$ (presión o altitud) y sobre el \textbf{bus I2C}.

\subsection{Muestreo y teorema de Nyquist}
Cada \code{read\_all()} del BMP180 realiza dos transacciones I2C y espera la conversión
sigma--delta: \SI{8}{\milli\second} (oss=1) a \SI{26}{\milli\second} (oss=3). La
frecuencia de muestreo práctica es de 5--30 Hz. El teorema de Nyquist--Shannon exige
$F_s \ge 2 f_\text{max}$ para reconstruir una señal de banda limitada a $f_\text{max}$.
Un componente de frecuencia $f > F_s/2$ se \emph{pliega} (aliasing) a
$f_a = \lvert ((f + F_s/2)\bmod F_s) - F_s/2\rvert$. El firmware planifica las lecturas
con \code{ticks\_us} y reporta el $F_s$ efectivo y su \emph{jitter}; si se solicita más
frecuencia de la que el sensor puede entregar, el error de $F_s$ lo evidencia.

\subsection{Caracterización del ruido y resolución efectiva}
El ruido de presión del BMP180 (RMS 3--6 Pa según \code{oss}) se traduce en ruido de
altitud a razón de $\approx\SI{8.43}{\meter\per\hecto\pascal}$. Se caracteriza mediante:
\begin{itemize}
  \item \textbf{Desviación estándar} $\sigma_h$ y RMS de un bloque estático.
  \item \textbf{Resolución efectiva} $\text{ENOB}_\text{dc} =
        \log_2\!\bigl(\text{rango}/(\sqrt{12}\,\sigma_h)\bigr)$.
  \item \textbf{Densidad espectral de potencia} (método de Welch): forma del piso de
        ruido (blanco vs $1/f$).
  \item \textbf{Desviación de Allan} $\sigma_A(\tau)$ frente al tiempo de promediado
        $\tau$: para ruido blanco $\sigma_A \propto \tau^{-1/2}$; el mínimo marca el
        $\tau$ óptimo antes de que domine la deriva térmica/meteorológica.
\end{itemize}

\subsection{Filtrado digital}
\begin{description}
  \item[Media móvil de $N$] FIR de coeficientes iguales; función de transferencia
        $H(z) = \frac{1}{N}\frac{1-z^{-N}}{1-z^{-1}}$, con ceros en $k\,F_s/N$ y retardo
        de grupo $(N-1)/2$ muestras.
  \item[Mediana de $N$] filtro no lineal; elimina \emph{spikes} de $\le\lfloor N/2\rfloor$
        muestras sin suavizar los flancos reales.
  \item[Exponencial (EMA)] IIR de 1\textsuperscript{er} orden $y[n]=y[n-1]+\alpha(x[n]-y[n-1])$;
        pasa-bajos con $f_c \approx \frac{F_s\,\alpha}{2\pi(1-\alpha)}$.
\end{description}
En la PC se compara el efecto en el dominio del tiempo (suavizado vs retardo), en el
dominio de la frecuencia (Bode del filtro) y ante un escalón real (tiempo de subida,
sobre-oscilación).

\subsection{El bus I2C en el osciloscopio}
Con CH1 en SCL y CH2 en SDA (base \SI{10}{\micro\second\per div}, trigger flanco de
bajada) se mide: la frecuencia de SCL ($\approx\SI{100}{\kilo\hertz}$), la duración de
una transacción de lectura completa, y el periodo entre transacciones (igual a $1/F_s$
del modo activo). El ``tiempo muerto'' del bus indica el margen para aumentar $F_s$.

% ============================================================================
\section{Alternativa de Trabajo en Casa (Multímetro + Simulación)}
\label{sec:casa}

Esta práctica es, de las ocho, la que \textbf{menos depende} del laboratorio presencial:
el BMP180 es un sensor digital y ninguno de sus modos (incluidos los de análisis de
señales, \S8) requiere osciloscopio ni generador de funciones — toda la adquisición y
el análisis ocurren en el propio microcontrolador y en la PC. Aun así, se documenta
aquí qué hacer si no se dispone del módulo físico.

\paragraph{Con multímetro (en casa).} Con el módulo BMP180 alimentado: verificar
\SI{3.3}{\volt} en VCC. Con el módulo \emph{desconectado} de alimentación: medir con el
óhmetro la resistencia entre SDA/SCL y 3V3 — debe rondar los \SI{4.7}{\kilo\ohm} de los
pull-ups del módulo GY-68 (\S3.3); un valor muy distinto indica un pull-up dañado o un
corto. Verificar continuidad de los 4 cables Dupont.

\paragraph{Con simulación.} No existe un modelo de BMP180 en las versiones gratuitas de
NI Multisim/Proteus/Wokwi, así que la ``simulación'' de esta práctica es
\textbf{numérica}: reproducir el algoritmo de compensación (\S3.5–3.6) en una hoja de
cálculo, Python o el propio modo 2 del firmware (que ya expone cada paso intermedio),
usando los coeficientes de ejemplo de Bosch (\S6.2) para verificar que la implementación
es correcta antes de tener el sensor físico en la mano.

\paragraph{Requiere el sensor físico (no sustituible).} La comunicación I2C real y las
mediciones de altitud/presión de las \S7.3--7.8 solo pueden hacerse con el módulo
BMP180 conectado; no hay atajo de simulación para esa parte.

\begin{table}[H]\centering
\caption{Qué requiere laboratorio presencial en esta práctica.}
\begin{tabular}{lccc}
\toprule
Actividad & Multímetro (casa) & Simulación & Requiere BMP180 físico \\
\midrule
Verificar VCC / pull-ups & Sí & --- & --- \\
Verificar algoritmo de compensación & --- & Sí (numérica) & --- \\
Casos 1--5 (lecturas reales) & --- & --- & Sí \\
Casos 6--7 (análisis de señales) & --- & --- & Sí (pero sin osciloscopio) \\
Bus I2C en osciloscopio (\S8.4) & --- & --- & Sí, y osciloscopio \\
\bottomrule
\end{tabular}
\end{table}

% ============================================================================
\section{Registro de Datos y Análisis}

\begin{table}[H]\centering
\caption{Coeficientes de calibración medidos.}
\label{tab:coef}
\begin{tabular}{lcc | lcc}
\toprule
Coef.\ & Valor leído & Tipo & Coef.\ & Valor leído & Tipo \\
\midrule
AC1 & & int16  & B1 & & int16 \\
AC2 & & int16  & B2 & & int16 \\
AC3 & & int16  & MB & & int16 \\
AC4 & & uint16 & MC & & int16 \\
AC5 & & uint16 & MD & & int16 \\
AC6 & & uint16 & & & \\
\bottomrule
\end{tabular}
\end{table}

\begin{table}[H]\centering
\caption{Verificación paso a paso del algoritmo de compensación.}
\begin{tabular}{lll}
\toprule
Variable & Ecuación & Valor calculado / medido \\
\midrule
UT & lectura directa & \\
$X_1$ (temp) & $(\text{UT}-\text{AC6})\,\text{AC5}/2^{15}$ & \\
$X_2$ (temp) & $\text{MC}\cdot 2^{11}/(X_1+\text{MD})$ & \\
$B_5$ & $X_1+X_2$ & \\
$T$ & $(B_5+8)/2^4$ & \\
UP & lectura directa & \\
$P$ (final) & algoritmo completo & \\
$h$ (altitud) & ec.~\eqref{eq:isa} & \\
\bottomrule
\end{tabular}
\end{table}

\begin{table}[H]\centering
\caption{Análisis de muestreo (Caso 6).}
\begin{tabular}{ccccc}
\toprule
$F_s$ solicitada (Hz) & $F_s$ real (Hz) & error (\%) & jitter $\sigma$ (\si{\micro\second}) & dropouts \\
\midrule
5 & & & & \\
10 & & & & \\
20 & & & & \\
\bottomrule
\end{tabular}
\end{table}

\begin{table}[H]\centering
\caption{Ruido de altitud vs modo de sobremuestreo.}
\begin{tabular}{ccccc}
\toprule
OSS & t/lectura (ms) & $\sigma_h$ (m) & RMS ruido (m) & resolución efectiva (bits/\SI{100}{\meter}) \\
\midrule
0 & & & & \\
1 & & & & \\
2 & & & & \\
3 & & & & \\
\bottomrule
\end{tabular}
\end{table}

\begin{table}[H]\centering
\caption{Filtros digitales (respuesta a un escalón de $\sim\SI{3}{\meter}$).}
\begin{tabular}{lccc}
\toprule
Filtro & t.\ subida 10--90\,\% (s) & sobre-oscilación (\%) & $\sigma$ residual en reposo (m) \\
\midrule
crudo & --- & --- & \\
media móvil $N=8$ & & & \\
mediana $N=5$ & & & \\
EMA $\alpha=0.2$ & & & \\
\bottomrule
\end{tabular}
\end{table}

\subsection{Gráficas requeridas}
\begin{enumerate}
  \item Altitud vs tiempo de un recorrido vertical (crudo + filtrado superpuestos).
  \item Presión vs altitud de la Tabla de comparativa, con la curva teórica ISA superpuesta.
  \item PSD del ruido de altitud en reposo (log--log) para oss=1 y oss=3.
  \item Desviación de Allan vs $\tau$ (identificar el mínimo).
  \item Bode (magnitud) de la media móvil $N=8$ y de la EMA $\alpha=0.2$, marcando $f_{-3\,\text{dB}}$.
  \item Captura de osciloscopio de una transacción I2C con anotaciones.
\end{enumerate}

% ============================================================================
\section{Análisis de Resultados}
\begin{enumerate}
  \item Comparar los coeficientes leídos con los rangos típicos de Bosch. ¿Alguno parece
        anómalo?
  \item Verificar que el algoritmo produce valores razonables ($T\sim$ 20--30 \si{\celsius},
        $P\sim$ 900--1050 hPa según elevación).
  \item Evaluar la repetibilidad: ¿la desviación estándar de cada punto de la comparativa
        es menor a \SI{0.5}{\meter}?
  \item Calcular la altura real de cada piso del edificio a partir de los $\Delta h$
        medidos. ¿Son consistentes (cada piso $\approx$ 3.0--3.5 m)?
  \item A partir de la PSD del Caso 6: ¿el ruido de altitud es blanco o tiene componente
        $1/f$? ¿Cambia con \code{oss}?
  \item ¿A partir de qué $\tau$ deja de mejorar la desviación de Allan? ¿Qué lo limita?
  \item Comparar la resolución teórica (\SI{0.01}{\hecto\pascal} $\approx$ \SI{8.5}{\centi\meter})
        con la resolución efectiva observada.
\end{enumerate}

% ============================================================================
\section{Preguntas de Discusión y Refuerzo}
\begin{enumerate}
  \item \textbf{Nyquist en un sensor lento}: a oss=1 el sensor tarda \SI{8}{\milli\second}
        por lectura. ¿Cuál es la frecuencia máxima de una oscilación de presión (p.ej.\ una
        puerta batiente) que podrías medir sin aliasing? ¿Qué pasa si intentas muestrear
        a \SI{20}{\hertz}?
  \item \textbf{Aritmética entera}: ¿por qué el \emph{datasheet} especifica el algoritmo
        en enteros de 32 bits y no en punto flotante? Discutir portabilidad y determinismo.
  \item \textbf{Coeficiente $B_5$}: ¿por qué aparece tanto en el cálculo de temperatura
        como en el de presión? ¿Qué ocurre si se lee la presión sin actualizar $B_5$?
  \item \textbf{QNH variable}: si el QNH cambia \SI{5}{\hecto\pascal} entre la calibración
        y una medición posterior, ¿cuánto error de altitud introduce? Calcular con la
        ec.~\eqref{eq:isa}.
  \item \textbf{Promediado}: promediar 16 lecturas reduce el ruido $\sqrt{16}=4\times$.
        ¿Por qué la desviación de Allan deja de bajar a partir de cierto $\tau$?
  \item \textbf{Media móvil}: la media móvil de $N=8$ a $F_s=\SI{5}{\hertz}$, ¿qué
        frecuencias elimina por completo? ¿Qué retardo introduce en segundos?
  \item \textbf{Mediana vs media}: ¿en qué situación la mediana supera claramente a la
        media móvil? ¿Qué desventaja tiene ante un escalón real?
  \item \textbf{Resolución vs precisión}: el BMP180 tiene resolución de
        \SI{0.01}{\hecto\pascal} pero precisión de $\pm\SI{1.0}{\hecto\pascal}$. ¿Qué
        significa esta diferencia en la práctica?
  \item \textbf{ARINC 429}: ¿cómo se codificaría una altitud de \SI{1250}{ft} en formato
        BNR dentro de una palabra Label 203? Especificar los bits de dato, SSM y paridad.
\end{enumerate}

% ============================================================================
\section{Entregables (Formato AIAA)}
\begin{enumerate}
  \item \textbf{Código fuente}: \code{main.py} + \code{lib/bmp180.py} + \code{lib/siglab.py}
        (o \code{p4.cpp} + \code{common/siglab.h}) con comentarios claros.
  \item \textbf{Tabla de coeficientes}: los 11 valores leídos de tu sensor.
  \item \textbf{Verificación del algoritmo}: tabla con valores intermedios a mano vs medidos.
  \item \textbf{Mediciones de altura}: $\ge 3$ puntos con $\sigma$ y $\Delta h$.
  \item \textbf{Análisis de señales}: tablas de muestreo, ruido vs OSS y filtros (Caso 6--7).
  \item \textbf{Gráficas}: las 6 de la \S9.4.
  \item \textbf{Análisis}: respuestas a $\ge 4$ preguntas de la \S11.
  \item \textbf{Conclusiones}: hallazgos, fuentes de error, sugerencias de mejora.
  \item \textbf{Referencias}: $\ge 3$ fuentes consultadas.
\end{enumerate}

% ============================================================================
\section{Referencias}
\begin{enumerate}
  \item Bosch Sensortec. (2013). \emph{BMP180 Digital Pressure Sensor Datasheet},
        BST-BMP180-DS000-09, Rev 2.5.
  \item ICAO. (1993). \emph{Manual of the ICAO Standard Atmosphere}, Doc 7488/3, 3rd ed.
  \item ARINC. (2004). \emph{ARINC Specification 429 Part 1-17: Mark 33 DITS}.
  \item NXP Semiconductors. (2021). \emph{UM10204: I2C-bus Specification and User Manual},
        Rev 7.0.
  \item MicroPython Project. \emph{MicroPython Documentation --- machine.I2C}.
  \item Oppenheim, A.V. \& Schafer, R.W. \emph{Discrete-Time Signal Processing}, 3rd ed.
  \item IEEE. \emph{IEEE Std 1241-2010 --- Standard for Terminology and Test Methods for
        Analog-to-Digital Converters}.
  \item Moir, I. \& Seabridge, A. (2008). \emph{Aircraft Systems}, 3rd ed., Wiley.
\end{enumerate}

% ============================================================================
\appendix
\section{Código Fuente del Repositorio}
El código completo está en \url{https://github.com/EU97/SISELA-Init}.
\begin{table}[H]\centering
\caption{Archivos relevantes de la Práctica 4.}
\begin{tabular}{ll}
\toprule
Ruta & Descripción \\
\midrule
\code{MicroPython/\{ESP32,RP2040\}/P4/main.py} & Programa principal (7 modos) \\
\code{.../P4/lib/bmp180.py} & Driver BMP180 (compensación completa) \\
\code{.../P4/lib/siglab.py} & Utilidades de análisis de señales en la placa \\
\code{.../P4/docs/analisis\_senales.md} & Guía del módulo de análisis \\
\code{MicroPython/ESP32/P4/tools/altimeter\_gui.py} & Altímetro visual (tkinter) \\
\code{C++/SISELA-CPP/src/practices/p4.cpp} & Implementación C++ (Wire, 6 modos) \\
\code{C++/SISELA-CPP/src/common/siglab.h} & \code{siglab} para C++ (header-only) \\
\code{tools/sisela\_signal/} & Toolkit PC de análisis de señales \\
\bottomrule
\end{tabular}
\end{table}

\section{Contexto Aeronáutico Complementario}
\textbf{El altímetro en la cabina de vuelo.} El altímetro barométrico es un instrumento
del grupo pitot-estático. La ventana de Kollsman permite al piloto ajustar el QNH.
\textbf{RVSM} (\emph{Reduced Vertical Separation Minimum}): entre FL290 y FL410 la
separación vertical mínima se reduce a \SI{1000}{ft}, lo que exige altímetros con una
tolerancia muy estricta y por eso el énfasis en la caracterización del ruido y la
resolución efectiva de esta práctica.

\end{document}
